% GENERATED FILE. Edit canonical lessons, then run npm run generate:books.
% canonical-source-hash: fnv1a64:e68af947b41cf9a4
% canonical-lessons: TE-C129-kitikitera, TE-C129-catti, TE-C129-binde, TE-C129-biruva, TE-C129-tebul, TE-R129-first-pass-food-and-nature, TE-R129-second-pass-nature-and-animals

\chapter{Window curtain, Clay pot, Water pot, Cupboard, Table}
\label{ch:te-window-curtain-clay-pot-water-pot-cupboard-table}
\hlchaptermodality{129}

\begin{chapteropening}
\textbf{By the end of this chapter:} \emph{I can say a window curtain, a clay pot, a water pot, a cupboard and a table.}

Complete the last lesson of chapter 129: I can say a window curtain, a clay pot, a water pot, a cupboard and a table.
\end{chapteropening}

\section[\texorpdfstring{\te{కిటికీతెర} (kiṭikītera)}{కిటికీతెర (kiṭikītera)}]{\te{కిటికీతెర} (kiṭikītera) — a window curtain}

\label{lesson:TE-C129-kitikitera}



\begin{quote}
Before the new one: say the Telugu for the skin, then the Telugu for blood.
\end{quote}

\subsection*{You'll want to know: \te{కిటికీతెర}}
\textbf{\te{కిటికీతెర}} — \emph{kiṭikītera} — \textquotedblleft{}a window curtain\textquotedblright{}.

\begin{grammarlens}[title={What you've built}]
One more word for this chapter.
\end{grammarlens}

\subsection*{Your turn}
\begin{itemize}
\raggedright
  \item \emph{Say it:} \emph{kiṭikītera}
  \item \emph{Say it:} \emph{kiṭikītera}, once more
  \item \emph{Say it:} say \emph{kiṭikītera}
  \item \emph{Recall:} say the Telugu for the skin, then the Telugu for blood, then say \emph{kiṭikītera} again
\end{itemize}

\begin{tcolorbox}[breakable,colback=teal!4,colframe=teal!35!black,title={Before you move on}]
What does \te{కిటికీతెర} mean? (\textquotedblleft{}A window curtain\textquotedblright{}.) Say it once more.
\end{tcolorbox}
\section[\texorpdfstring{\te{చట్టి} (caṭṭi)}{చట్టి (caṭṭi)}]{\te{చట్టి} (caṭṭi) — a clay pot}

\label{lesson:TE-C129-catti}



\begin{quote}
Before the new one: say the Telugu for blood, then the Telugu for a window curtain.
\end{quote}

\subsection*{You'll want to know: \te{చట్టి}}
\textbf{\te{చట్టి}} — \emph{caṭṭi} — \textquotedblleft{}a clay pot\textquotedblright{}.

\begin{grammarlens}[title={What you've built}]
One more word for this chapter.
\end{grammarlens}

\subsection*{Your turn}
\begin{itemize}
\raggedright
  \item \emph{Say it:} \emph{caṭṭi}
  \item \emph{Say it:} \emph{caṭṭi}, once more
  \item \emph{Say it:} say \emph{caṭṭi}
  \item \emph{Recall:} say the Telugu for blood, then the Telugu for a window curtain, then say \emph{caṭṭi} again
\end{itemize}

\begin{tcolorbox}[breakable,colback=teal!4,colframe=teal!35!black,title={Before you move on}]
What does \te{చట్టి} mean? (\textquotedblleft{}A clay pot\textquotedblright{}.) Say it once more.
\end{tcolorbox}
\section[\texorpdfstring{\te{బిందె} (bindē)}{బిందె (bindē)}]{\te{బిందె} (bindē) — a water pot}

\label{lesson:TE-C129-binde}



\begin{quote}
Before the new one: say the Telugu for a window curtain, then the Telugu for a clay pot.
\end{quote}

\subsection*{You'll want to know: \te{బిందె}}
\textbf{\te{బిందె}} — \emph{bindē} — \textquotedblleft{}a water pot\textquotedblright{}.

\begin{grammarlens}[title={What you've built}]
One more word for this chapter.
\end{grammarlens}

\subsection*{Your turn}
\begin{itemize}
\raggedright
  \item \emph{Say it:} \emph{bindē}
  \item \emph{Say it:} \emph{bindē}, once more
  \item \emph{Say it:} say \emph{bindē}
  \item \emph{Recall:} say the Telugu for a window curtain, then the Telugu for a clay pot, then say \emph{bindē} again
\end{itemize}

\begin{tcolorbox}[breakable,colback=teal!4,colframe=teal!35!black,title={Before you move on}]
What does \te{బిందె} mean? (\textquotedblleft{}A water pot\textquotedblright{}.) Say it once more.
\end{tcolorbox}
\section[\texorpdfstring{\te{బీరువా} (bīruvā)}{బీరువా (bīruvā)}]{\te{బీరువా} (bīruvā) — a cupboard}

\label{lesson:TE-C129-biruva}



\begin{quote}
Before the new one: say the Telugu for a clay pot, then the Telugu for a water pot.
\end{quote}

\subsection*{You'll want to know: \te{బీరువా}}
\textbf{\te{బీరువా}} — \emph{bīruvā} — \textquotedblleft{}a cupboard\textquotedblright{}.

\begin{grammarlens}[title={What you've built}]
One more word for this chapter.
\end{grammarlens}

\subsection*{Your turn}
\begin{itemize}
\raggedright
  \item \emph{Say it:} \emph{bīruvā}
  \item \emph{Say it:} \emph{bīruvā}, once more
  \item \emph{Say it:} say \emph{bīruvā}
  \item \emph{Recall:} say the Telugu for a clay pot, then the Telugu for a water pot, then say \emph{bīruvā} again
\end{itemize}

\begin{tcolorbox}[breakable,colback=teal!4,colframe=teal!35!black,title={Before you move on}]
What does \te{బీరువా} mean? (\textquotedblleft{}A cupboard\textquotedblright{}.) Say it once more.
\end{tcolorbox}
\section[\texorpdfstring{\te{టేబుల్} (ṭēbul)}{టేబుల్ (ṭēbul)}]{\te{టేబుల్} (ṭēbul) — a table}

\label{lesson:TE-C129-tebul}



\begin{quote}
Before the new one: say the Telugu for a water pot, then the Telugu for a cupboard.
\end{quote}

\subsection*{You'll want to know: \te{టేబుల్}}
\textbf{\te{టేబుల్}} — \emph{ṭēbul} — \textquotedblleft{}a table\textquotedblright{}.

\begin{grammarlens}[title={What you've built}]
One more word for this chapter.
\end{grammarlens}

\subsection*{Your turn}
\begin{itemize}
\raggedright
  \item \emph{Say it:} \emph{ṭēbul}
  \item \emph{Say it:} \emph{ṭēbul}, once more
  \item \emph{Say it:} say \emph{ṭēbul}
  \item \emph{Recall:} say the Telugu for a water pot, then the Telugu for a cupboard, then say \emph{ṭēbul} again
\end{itemize}

\begin{tcolorbox}[breakable,colback=teal!4,colframe=teal!35!black,title={Before you move on}]
What does \te{టేబుల్} mean? (\textquotedblleft{}A table\textquotedblright{}.) Say it once more.
\end{tcolorbox}
\section[(dialogue)]{Food and nature — a first pass}

\label{lesson:TE-R129-first-pass-food-and-nature}



\begin{quote}
Say the words for a cupboard and a table.
\end{quote}

\subsection*{Your turn}
Say these aloud:

\begin{itemize}
\raggedright
  \item the sea, a lake, a forest, a leaf, soil
  \item thunder, lightning, a storm, a wave, a hill
  \item a waterfall, an island, a cow, a horse, an elephant
  \item a tiger, a lion, a bear, a monkey, a rabbit
  \item a deer, a fox, a snake, a tortoise, a parrot
\end{itemize}

\begin{tcolorbox}[breakable,colback=teal!4,colframe=teal!35!black,title={Before you move on}]
The sea? (\textbf{\te{సముద్రం}}, \emph{samudraṁ}.) A lake? (\textbf{\te{సరస్సు}}, \emph{sarassu}.) A forest? (\textbf{\te{అడవి}}, \emph{aḍavi}.) A leaf? (\textbf{\te{ఆకు}}, \emph{āku}.) Soil? (\textbf{\te{మట్టి}}, \emph{maṭṭi}.) Thunder? (\textbf{\te{ఉరుము}}, \emph{urumu}.) Lightning? (\textbf{\te{మెరుపు}}, \emph{merupu}.) A storm? (\textbf{\te{తుపాను}}, \emph{tupānu}.) A wave? (\textbf{\te{అల}}, \emph{ala}.) A hill? (\textbf{\te{కొండ}}, \emph{koṇḍa}.) A waterfall? (\textbf{\te{జలపాతం}}, \emph{jalapātaṁ}.) An island? (\textbf{\te{దీవి}}, \emph{dīvi}.) A cow? (\textbf{\te{ఆవు}}, \emph{āvu}.) A horse? (\textbf{\te{గుర్రం}}, \emph{gurraṁ}.) An elephant? (\textbf{\te{ఏనుగు}}, \emph{ēnugu}.) A tiger? (\textbf{\te{పులి}}, \emph{puli}.) A lion? (\textbf{\te{సింహం}}, \emph{siṁhaṁ}.) A bear? (\textbf{\te{ఎలుగుబంటి}}, \emph{elugubaṇṭi}.) A monkey? (\textbf{\te{కోతి}}, \emph{kōti}.) A rabbit? (\textbf{\te{కుందేలు}}, \emph{kundēlu}.) A deer? (\textbf{\te{జింక}}, \emph{jiṅka}.) A fox? (\textbf{\te{నక్క}}, \emph{nakka}.) A snake? (\textbf{\te{పాము}}, \emph{pāmu}.) A tortoise? (\textbf{\te{తాబేలు}}, \emph{tābēlu}.) A parrot? (\textbf{\te{చిలక}}, \emph{cilaka}.)
\end{tcolorbox}
\section[(dialogue)]{Nature and animals — a second pass}

\label{lesson:TE-R129-second-pass-nature-and-animals}



\begin{quote}
Say the words for a pigeon and an owl.
\end{quote}

\subsection*{Your turn}
Say these aloud:

\begin{itemize}
\raggedright
  \item a pigeon, an owl, a peacock, a butterfly, a crow
  \item a duck, a pig, a donkey, a camel, the face
  \item the tongue, a knee, a shoulder, an eyebrow, a cheek
  \item the chin; a beard, a wrist, an elbow, the skin, blood
  \item a window curtain, a clay pot, a water pot, a cupboard, a table
\end{itemize}

\begin{tcolorbox}[breakable,colback=teal!4,colframe=teal!35!black,title={Before you move on}]
A pigeon? (\textbf{\te{పావురం}}, \emph{pāvuraṁ}.) An owl? (\textbf{\te{గుడ్లగూబ}}, \emph{guḍlagūba}.) A peacock? (\textbf{\te{నెమలి}}, \emph{nemali}.) A butterfly? (\textbf{\te{సీతాకోకచిలుక}}, \emph{sītākōkacilaka}.) A crow? (\textbf{\te{కాకి}}, \emph{kāki}.) A duck? (\textbf{\te{బాతు}}, \emph{bātu}.) A pig? (\textbf{\te{పంది}}, \emph{pandi}.) A donkey? (\textbf{\te{గాడిద}}, \emph{gāḍida}.) A camel? (\textbf{\te{ఒంటె}}, \emph{oṇṭe}.) The face? (\textbf{\te{ముఖం}}, \emph{mukhaṁ}.) The tongue? (\textbf{\te{నాలుక}}, \emph{nāluka}.) A knee? (\textbf{\te{మోకాలు}}, \emph{mōkālu}.) A shoulder? (\textbf{\te{భుజం}}, \emph{bhujaṁ}.) An eyebrow? (\textbf{\te{కనుబొమ్మ}}, \emph{kanubomma}.) A cheek? (\textbf{\te{బుగ్గ}}, \emph{bugga}.) The chin; a beard? (\textbf{\te{గడ్డం}}, \emph{gaḍḍaṁ}.) A wrist? (\textbf{\te{మణికట్టు}}, \emph{maṇikaṭṭu}.) An elbow? (\textbf{\te{మోచేయి}}, \emph{mōcēyi}.) The skin? (\textbf{\te{చర్మం}}, \emph{carmaṁ}.) Blood? (\textbf{\te{రక్తం}}, \emph{raktaṁ}.) A window curtain? (\textbf{\te{కిటికీతెర}}, \emph{kiṭikītera}.) A clay pot? (\textbf{\te{చట్టి}}, \emph{caṭṭi}.) A water pot? (\textbf{\te{బిందె}}, \emph{bindē}.) A cupboard? (\textbf{\te{బీరువా}}, \emph{bīruvā}.) A table? (\textbf{\te{టేబుల్}}, \emph{ṭēbul}.)
\end{tcolorbox}
